%
%============= CHAPTER BREAK =================%
%
\chapter{Charax of Pergamum (ca 150--200 CE)}%
\vspace*{3mm}%
\justifying%
%
\noindent{}According to the most recent edition of the \textit{Encyclopaedia Judaica},\endnote{%
4:568, s.v., “Charax of Pergamum”.} %
Charax wrote sometime in the 2nd half of the 2nd century. \\[1.5mm] %

\noindent{}\textbf{μεμβράνα} / \textit{membrana} / “membrane” may have appeared first %
in a writing by Charax, but it is only attested second-hand almost a %
millennium later. %
According to a fragmentary and paraphrasing summary %
of his writing as preserved by Eustathius of Thessalonica (late 12th cent.): % 

\vspace*{3mm} 
 \begin{quote}Ὁ Χάραξ τὸ χρυσοῦν δέρμα μέθοδον εἶναι λέγει %
 χρυσογραφίας \textbf{μεμβράναις} ἐμπεριειλημμένην, %
 δι’ ἣν, ὡς λόγου ἀξίαν, %
 τὸν τῆς Ἀργοῦς καταρτισθῆναι στόλον φησί. (\textit{Frag}. 14) \\[1mm] %
 Charax says the craft of golden skin was gold-lettering gilded \textbf{on membranes}, %
 on account of which, as is worthy of note, %
 he says the fleet of Argos was equipped. \end{quote} % 
 \vspace*{3mm} 

\noindent{}This lemma only appears once in the NT, %
in a Pastoral letter (2Tim 4:13): 

\vspace*{3mm} 
 \begin{quote}τὸν φαιλόνην ὃν ἀπέλιπον ἐν Τρῳάδι παρὰ Κάρπῳ
 ἐρχόμενος φέρε, καὶ τὰ βιβλία, μάλιστα τὰς \textbf{μεμβράνας}. \\[1mm] %
 Bring the cloak which I left behind in Troas with Carpus, %
 carrying it and the little books, especially the \textbf{membranes}. \end{quote} % 
 \vspace*{3mm}

\noindent{}It also appears twice in quick succession in the apocryphal %
\textit{Acts of John by Prochorus}.\endnote{
Zahn 1880:156.}
While the evidence is admittedly minimal and questionable, %
it suggests that early-Orthodox portraiture of Paul %
was modeled, whether in passing and/or as satire, %
on the prestigious, opulent scribal craft detailed by Charax. %